\documentclass[aspectratio=169,11pt]{ctexbeamer}

\usetheme{Madrid}
\usecolortheme{default}
\usefonttheme{professionalfonts}
\setbeamertemplate{navigation symbols}{}
\setbeamertemplate{footline}[frame number]
\pdfstringdefDisableCommands{\def\texttt#1{#1}}

\definecolor{HustBlue}{RGB}{17,64,124}
\definecolor{HustLight}{RGB}{233,241,250}
\definecolor{HustAccent}{RGB}{0,123,167}

\setbeamercolor{palette primary}{bg=HustBlue,fg=white}
\setbeamercolor{palette secondary}{bg=HustBlue,fg=white}
\setbeamercolor{palette tertiary}{bg=HustBlue,fg=white}
\setbeamercolor{structure}{fg=HustBlue}
\setbeamercolor{frametitle}{bg=HustLight,fg=HustBlue}
\setbeamercolor{title}{fg=HustBlue}
\setbeamercolor{block title}{bg=HustBlue,fg=white}
\setbeamercolor{block body}{bg=HustLight,fg=black}

\title{第一个里程碑阶段进度汇报}
\subtitle{面向国产硬件极致性能优化的推理引擎}
\author{张书豪}
\date{2026-04-01}

\begin{document}

\begin{frame}
  \titlepage
\end{frame}

\begin{frame}{里程碑定位与阶段目标}
  \begin{itemize}
    \item 当前第一个里程碑的定位，不是 9 个优化点的全面交付。
    \item 当前阶段重点完成四项工作：
      \begin{itemize}
        \item 公共基线建立
        \item 统一指标口径收敛
        \item 统一评测与结果留痕方式确定
        \item 阶段材料和证据链归口
      \end{itemize}
    \item 阶段目标是先把底层推理底座、统一接入入口、真实业务测试对象和基础观测能力贯通起来。
  \end{itemize}
\end{frame}

\begin{frame}{阶段性总体判断}
  \begin{block}{当前总判断}
    项目总体实施方案与技术架构已经完成阶段性定型，\texttt{vllm-hust} 统一推理底座、\texttt{workstation} 本地控制面与 \texttt{EvoScientist} 首个真实多智能体测试对象三层关系已经明确，底层 serving 能力和部分上层接入能力已经具备。
  \end{block}

  \vspace{0.5em}
  \begin{itemize}
    \item 当前工作重点已从“论证底层是否可行”转向“补齐真实场景验证、统一指标与画像 schema、跨硬件实测和流式数据闭环等证据链”。
    \item 当前最适合的阶段表述是：架构与链路已贯通，数据与证据仍需继续补齐。
  \end{itemize}
\end{frame}

\begin{frame}{当前已完成的阶段性工作}
  \begin{enumerate}
    \item 总体实施方案与技术架构完成阶段性定型。
    \item 底层 serving 能力已具备基础可用性：多模态、流式输出、reasoning、工具调用、结构化输出、benchmark 与 metrics 入口均已具备。
    \item 首个真实多智能体业务对象已接入：\texttt{EvoScientist} 已接入本地 \texttt{vllm-hust}，默认桥接链路已打通。
    \item 统一后端与环境封装需求分析已完成：OpenAI-compatible 统一入口与 plugin / platform 后端边界已经明确。
  \end{enumerate}
\end{frame}

\begin{frame}{当前仓库与平台分工已经明确}
  \begin{itemize}
    \item \texttt{vllm-hust}：统一推理引擎仓库，承载 serving 核心能力。
    \item \texttt{vllm-hust-workstation}：应用平台仓库，承载本地工作台、任务入口、指标聚合和演示控制。
    \item \texttt{vllm-hust-benchmark}：评测与分析仓库，承载 baseline、对比实验和结果沉淀。
    \item \texttt{vllm-ascend-hust}：国产算力适配仓库，承载昇腾等国产硬件路径上的运行与优化。
    \item \texttt{vllm-hust-docs}：文档与证据归档仓库，承载实施方案、阶段材料和证据链归档。
  \end{itemize}
\end{frame}

\begin{frame}{当前已具备基础的支撑能力}
  \begin{columns}[T]
    \begin{column}{0.48\textwidth}
      \begin{block}{国产算力方向}
        \begin{itemize}
          \item 已形成明确瓶颈分析框架
          \item 主要关注前处理、TTFT / ITL、KV 传输与缓存、动态 batch 与通信路径
          \item 为后续专项优化和跨硬件实测提供分析基础
        \end{itemize}
      \end{block}
    \end{column}
    \begin{column}{0.48\textwidth}
      \begin{block}{负载画像与流式交互}
        \begin{itemize}
          \item 已具备 TTFT、ITL、TPOT、token throughput 等基础抓手
          \item 已具备 error type、reasoning length、tool density 等真实 workload 采集入口
          \item 已具备继续形成统一画像报表的基础条件
        \end{itemize}
      \end{block}
    \end{column}
  \end{columns}
\end{frame}

\begin{frame}{当前已具备基础但仍需补证据的事项}
  \begin{itemize}
    \item \textbf{多智能体仿真环境}：已有真实业务测试对象和默认接入链路，仍需补更系统的任务集、统一 workload schema 和阶段报表。
    \item \textbf{流式强化学习方向}：serving 侧数据面和 rollout 方向已具备基础，仍需补完整训练闭环和标准化数据回放链路。
    \item \textbf{国产算力实测证据}：统一后端路线和瓶颈方向已明确，仍需补至少 1 类国产芯片上的首轮 baseline 数据和跨硬件对比证据。
    \item \textbf{统一画像与报表体系}：基础采集字段已具备，仍需补统一 schema、错误分类和阶段报表模板。
  \end{itemize}
\end{frame}

\begin{frame}{当前不宜过度表述的内容}
  \begin{columns}[T]
    \begin{column}{0.48\textwidth}
      \begin{alertblock}{当前不宜表述为}
        \begin{itemize}
          \item 通用多智能体仿真平台已经建成
          \item 完整流式 RL 训练闭环已经落地
          \item 跨集群云原生调度算法验证已经完成
          \item 跨硬件同任务性能结论已经充分完备
        \end{itemize}
      \end{alertblock}
    \end{column}
    \begin{column}{0.48\textwidth}
      \begin{block}{当前准确表述}
        \begin{itemize}
          \item 底层能力和真实业务接入基础已经具备
          \item 统一指标和验证证据链正在收口
          \item 下一阶段重点补真实场景、跨硬件实测和正式数据闭环
        \end{itemize}
      \end{block}
    \end{column}
  \end{columns}
\end{frame}

\begin{frame}{下一阶段推进重点}
  \begin{enumerate}
    \item 补齐统一指标与统一基线：固化 TTFT、TPOT、吞吐、MFU、单位 token 成本口径与结果归档格式。
    \item 补齐国产硬件首轮 baseline 数据：完成至少 1 类国产芯片上的 7B-13B 基本链路跑通，形成问题清单。
    \item 补齐真实 workload 验证和画像闭环：围绕 \texttt{EvoScientist} 形成更稳定的任务验证链路、统一 schema 和阶段报表。
    \item 建立后续优化结果统一回收机制：各优化方向按既定分工推进，结果按统一口径回收到里程碑体系中。
  \end{enumerate}
\end{frame}

\begin{frame}{阶段性结论}
  \begin{block}{结论}
    当前第一个里程碑已完成阶段性架构定型和基础链路贯通，统一推理底座、统一接入入口和首个真实业务测试对象已经形成闭环基础。下一阶段将在现有底座上继续补齐统一基线、国产硬件实测和真实 workload 验证证据，推动里程碑从“架构和链路已明确”进一步走向“数据和证据可归档”。
  \end{block}
\end{frame}

\end{document}